% Generated by tools/edn_latex.py from reports/edn.md.
% Do not edit: change reports/edn.md and run `make edn`.
\ednentry{
  id = {EDN-49},
  title = {No bias correction on the log-to-euro inverse transform},
  date = {2026-09-30},
  milestone = {M3: Quality Assurance},
  topic = {Modelling Approach, Evaluation Protocol},
  participants = {@lukas2510},
  decision = {\texttt{predict\_\allowbreak{}eur} is \texttt{exp} of the model's log-space prediction, with no Duan smearing factor and no lognormal variance correction. The target functional is the conditional \textbf{median} price, and that is what every metric this project reports measures.},
  rationale = {\emph{Alternatives considered:}
\begin{itemize}[leftmargin=*, topsep=2pt]
\item \textbf{Option A (chosen): naive \texttt{exp}.}
\newline Pros: it is the conditional median of price, which is exactly what MdAPE, the +/-20\,\% share and the conformal intervals are about; it needs no statistic fitted on the training residuals, so nothing extra has to be persisted, versioned or kept in step between \texttt{train}, \texttt{evaluate} and the API; and it is what the model card already told readers the model does.
\newline Cons: it is \emph{not} an unbiased estimate of the mean price, so anyone reading the estimate as an expected sale value is reading it wrong. The model card and the response schema say ``typical price'' for that reason.
\item \textbf{Option B: Duan's smearing estimator.} Multiply by the mean of \texttt{exp(residual)} over the training rows.
\newline Pros: the textbook correction for a log-linear model, non-parametric, and it is what a reviewer who knows the retransformation problem will ask about.
\newline Cons: it targets the conditional mean, so it shifts every prediction upward, away from the typical asking price. Measured on the real snapshot's 19,665 test rows it makes every variant worse: factor 1.0258 and MdAPE 9.52\,\% to 9.83\,\% for \texttt{b1}; 1.0052 and 6.83\,\% to 6.86\,\% for \texttt{lgbm-\allowbreak{}basic}; 1.0028 and 6.32\,\% to 6.37\,\% for \texttt{lgbm-\allowbreak{}extended}. It also adds a fitted constant to the artefact that has to travel with the model and be applied identically in three places.
\item \textbf{Option C: lognormal variance correction,} multiply by \texttt{exp(sigma\textasciicircum{}2 /\allowbreak{} 2)}.
\newline Pros: closed form, no residual pass needed.
\newline Cons: everything option B has, plus a distributional assumption the residuals do not have to satisfy; it is strictly weaker than smearing, which estimates the same quantity without the assumption.
\end{itemize}
\par
\emph{Why:} The correction answers a question nobody here asks. Fitting with squared error on \texttt{log(price)} estimates \texttt{E[log P | x]}, so \texttt{exp} of it is the geometric mean, which for the conditional distribution of a price is the median. Every number in problem-spec section 6 and every criterion in section 8 is median- or quantile-flavoured: MdAPE is the primary metric, SC-02 is a share within a band, and SC-05's intervals are conformal and therefore calibrated after the fact on held-out data, so they need no correction either. A mean-targeting correction would make each of those numbers worse, and the measurements say it does, on all three variants that have an inverse transform at all.
\par
\texttt{b0} is the reason the decision is visible in the code rather than implicit. The median baseline is fitted on \texttt{price} directly, not on \texttt{log\_\allowbreak{}price}, because a median commutes with a monotone transform: \texttt{exp(median(log price))} and \texttt{median(price)} are the same number, and the readable one is the one worth storing in a lookup table a person can audit. So B0 needs no inverse transform, B1 and the two LightGBM variants share one, and the seam puts it inside the model -- \texttt{predict\_\allowbreak{}eur} returns euros and \texttt{predict\_\allowbreak{}log\_\allowbreak{}price} is defined as its log -- so \texttt{evaluate} and the API cannot come to disagree about it.
\par
The honest limit of this entry: the effect is small. The largest degradation measured is 0.32 pp, on the interpretable baseline rather than on the deployed candidate. It is recorded because the retransformation problem is a standard question about a log-target model, the answer has numbers behind it, and the alternative would have quietly cost accuracy.},
  ai = {Information seeking, Alternative generation, Alternative assessment, Recommendation},
  response = {Accepted with modifications},
  assessment = {AI raised the retransformation question unprompted, named the two standard corrections and argued from the target functional rather than from convention, which is the argument that decides it. Its first write-up carried the effect sizes from a different run (0.5 to 2.2 pp) as if they were ours; those were re-measured on the real snapshot and are an order of magnitude smaller, so the entry now claims what this project measured and says that the effect is small. The direction of the effect -- worse on every variant -- reproduced.},
  aievidence = {Claude Code session on 2026-09-30 implementing issue \#37: the smearing factor and the paired MdAPE per variant were computed on the real snapshot in a scratchpad run of the whole chain, with the numbers reproduced in the model card.},
  evidence = {\href{https://github.com/mlops-2627q1-mds-upc/MLOPS_Recommenditos/blob/main/recommenditos/modeling/model.py}{\texttt{recommenditos/\allowbreak{}modeling/\allowbreak{}model.\allowbreak{}py}}, \texttt{Model.\allowbreak{}predict\_\allowbreak{}eur} and \texttt{MedianBaselineModel}; \href{https://github.com/mlops-2627q1-mds-upc/MLOPS_Recommenditos/blob/main/docs/docs/model-card.md}{model card}, Training Procedure, Training; \texttt{tests/\allowbreak{}test\_\allowbreak{}model.\allowbreak{}py::test\_\allowbreak{}predict\_\allowbreak{}log\_\allowbreak{}price\_\allowbreak{}is\_\allowbreak{}exactly\_\allowbreak{}the\_\allowbreak{}log\_\allowbreak{}of\_\allowbreak{}predict\_\allowbreak{}eur}; \href{https://github.com/mlops-2627q1-mds-upc/MLOPS_Recommenditos/issues/37}{issue \#37}.},
}
